\section{Smart Contracts}
\label{sec:contracts}
%% ============================================================

\Beluga{} ships with three core contracts, deployed at genesis. All contracts use the \texttt{@luxfi/standard} library \cite{lux-standard} for access control, upgradability (UUPS proxy), and reentrancy protection.

\subsection{ModelRegistry.sol}

An ERC-721 contract where each token represents a registered AI model.

\begin{lstlisting}[language=Solidity,caption=ModelRegistry core interface]
interface IModelRegistry {
    struct ModelMetadata {
        bytes32 weightsHash;     // SHA-256 of serialized weights
        string  architecture;    // e.g. "qwen3-8b-lora"
        uint256 paramCount;      // total parameter count
        uint256 pricePerRequest; // BLG wei per inference
        uint256 inferenceCount;  // lifetime request counter
        address host;            // current inference host
        bool    verified;        // governance-approved
    }

    function registerModel(
        bytes32 weightsHash,
        string calldata architecture,
        uint256 paramCount,
        uint256 pricePerRequest
    ) external returns (uint256 tokenId);

    function requestInference(
        uint256 tokenId,
        bytes calldata input
    ) external payable returns (bytes32 requestId);

    function getModel(uint256 tokenId)
        external view returns (ModelMetadata memory);
}
\end{lstlisting}

Key design decisions:
\begin{itemize}[nosep]
    \item The \texttt{weightsHash} is the canonical identifier---two models with the same weights hash are the same model regardless of who registered them (first-mover gets the NFT).
    \item \texttt{pricePerRequest} is a floor set by the owner; the actual price is computed by the inference pricing precompile based on demand.
    \item \texttt{inferenceCount} is incremented atomically by the precompile, not by the contract, preventing manipulation.
    \item Transfer of the NFT transfers model ownership, including future inference revenue.
\end{itemize}

\subsection{TrainingBounty.sol}

An escrow contract for training bounties using the \texttt{@luxfi/standard} ComputeMarket interface.

\begin{lstlisting}[language=Solidity,caption=TrainingBounty core interface]
interface ITrainingBounty {
    struct Bounty {
        address sponsor;
        uint256 reward;          // BLG escrowed
        bytes32 baseModelHash;   // starting model
        string  targetBenchmark; // e.g. "mmlu"
        uint256 targetDelta;     // basis points improvement
        uint256 deadline;        // block number
        BountyStatus status;
    }

    enum BountyStatus { Open, Claimed, Expired, Disputed }

    function postBounty(
        bytes32 baseModelHash,
        string calldata targetBenchmark,
        uint256 targetDelta,
        uint256 deadline
    ) external payable returns (uint256 bountyId);

    function submitResult(
        uint256 bountyId,
        bytes32 resultWeightsHash,
        uint256 measuredDelta,
        bytes calldata teeAttestation
    ) external;

    function claimReward(uint256 bountyId) external;
}
\end{lstlisting}

The dispute resolution mechanism uses a three-phase process:
\begin{enumerate}[nosep]
    \item \textbf{Challenge} (48h): Any party may challenge by staking 10\% of the bounty amount.
    \item \textbf{Re-evaluation} (72h): Three randomly selected \PoAI{} validators independently re-run the evaluation.
    \item \textbf{Settlement}: Majority vote determines outcome. Loser forfeits their stake.
\end{enumerate}

\subsection{InferenceRewards.sol}

Distributes \BLG{} from the inference pool to hosts based on verified service:

\begin{lstlisting}[language=Solidity,caption=InferenceRewards distribution]
interface IInferenceRewards {
    function claimHostReward(
        uint256 modelId,
        uint256 epoch
    ) external;

    function getEpochReward(uint256 epoch)
        external view returns (uint256);

    function getHostShare(
        address host,
        uint256 epoch
    ) external view returns (uint256);
}
\end{lstlisting}

Rewards are calculated per epoch (1 epoch = 7200 blocks $\approx$ 1 hour):

\begin{equation}
R_{\text{host}}^{(e)} = \frac{P_e}{N_e} \cdot \frac{Q_h}{\sum_{h' \in H} Q_{h'}}
\end{equation}

where $P_e$ is the pool emission for epoch $e$, $N_e$ is the number of active hosts, $Q_h$ is host $h$'s quality score (uptime $\times$ throughput $\times$ accuracy), and $H$ is the set of all active hosts in the epoch.

%% ============================================================
